\section{Introduction}
\label{sec:intro}

Algorithms have become the bureaucrats of the twenty-first century.
They adjudicate loans, screen job applications, triage medical care,
and rank electoral information. In function, they exercise what is
effectively \emph{sovereign power}: the authority to make binding
decisions that shape individual life chances at scale and at speed
exceeding any human institution. In structure, they resemble
\emph{absolutist monarchies}: opaque, centralised, and answerable
only to the entity that operates them.

The history of political thought offers a canonical solution to such
concentrated power. Montesquieu argued in \textit{De l'Esprit des
Lois}~\cite{montesquieu1748} that liberty can only be preserved if
the power to \emph{make laws} (Legislative), \emph{enforce laws}
(Executive), and \emph{interpret laws} (Judicial) are held by
distinct and independent bodies. Crucially, Montesquieu's argument
was not about the virtue of the rulers---it was about the
\emph{structure} of the institution. A virtuous king with absolute
power is still a monarch; a corrupt republic with separated powers
still has structural safeguards.

Current efforts to govern algorithmic systems treat the separation of
powers as a matter of \emph{procedural policy}: guidelines for
fairness, principles for transparency, and auditing regimes that
arrive after harm has occurred. We argue this is insufficient for
the same reason Montesquieu found Aristotle's reliance on
good rulers insufficient: procedural guidelines are easily bypassed,
ignored, or silently rewritten by the very entity they are meant to
constrain. The counter-examples of Section~\ref{sec:independence}
demonstrate that this is not theoretical---it is standard practice.

This paper proposes that Montesquieu's separation of powers can be
implemented as a \textbf{structural invariant} of the software
itself. We introduce \emph{Constitutional Code}: an architecture in
which:
\begin{itemize}
  \item \textbf{Legislative Power} is the set of linear resource
  types~\cite{btv2026} encoding compliance rules---immutable
  invariants that the operator cannot rewrite at runtime without
  producing a detectable binary substitution;
  \item \textbf{Executive Power} is the operator's runtime
  execution, type-constrained to produce non-repudiable evidence
  receipts~\cite{btv-t-repo} for every consequential decision; and
  \item \textbf{Judicial Power} is the verification layer, where
  independent auditors inspect the system's statistical behaviour
  via Zero-Knowledge proofs~\cite{btv-ar2026} without requiring the
  operator's cooperation.
\end{itemize}

We formalise this mapping in Section~\ref{sec:mapping}, prove the
mutual independence of the three powers in
Section~\ref{sec:independence}, and derive operational implications
for regulators in Section~\ref{sec:regulators}. This is not merely a
technical proposal. It is a \emph{change of paradigm}: from
institutions built on the operator's trust and retrospective audit
trails, to institutions built on protocol verification and
mathematical guarantees.

In the geocentric model of algorithmic governance, the operator is
the centre: the algorithm's truth is what the operator's database
says it is. Constitutional Code is the heliocentric correction:
the \emph{protocol} is the centre, and the operator is one actor
executing verifiable computations over a shared state whose integrity
no single party controls.

\paragraph{Contributions.}
\begin{enumerate}
  \item A formal definition of an Algorithmic Institution as a
  4-tuple $\langle \mathcal{L}, \mathcal{E}, \mathcal{J},
  \Sigma \rangle$ with a constitutional invariant
  (Definition~\ref{def:invariant}).
  \item Three independence theorems (Legislative, Judicial,
  Executive) and a Constitutional Completeness corollary proving
  that capability sets are mutually disjoint
  (Section~\ref{sec:independence}).
  \item Formal counter-examples showing that standard best-effort
  architectures fail each independence condition.
  \item An operational framework for regulators as
  ``Constituents'' defining the Constitutional Type Schema
  (Section~\ref{sec:regulators}).
  \item A statutory-to-type mapping table linking EU AI Act,
  GDPR, and LGPD requirements to concrete type invariants
  (Table~\ref{tab:law-to-types}).
\end{enumerate}
